% GDZ PPN236018647,native579–580upper;HKUST580–581.
\clearpage\SourceLeaf{579}\ArithmeticHead{550}{L}{TAGEBUCH MIT ERLÄUTERUNGEN.}
\NumberedItem{109}
\indent Inter duos numeros datos semper dantur infinite multi termini medii tum arithmetico geometrici tum harmonico geometrici, quorum nexum mutuum ex asse perspiciendi felicitas nobis est facta.
\par\noindent\hfill[1800] Iunio 3.\quad Brunov[ici]

\medskip\begingroup\small
\makeatletter\tagsleft@true\let\veqno\@@leqno\makeatother
\renewcommand{\thefootnote}{*)}
\indent Eine mit dieser Notiz gleichzeitige Aufzeichnung, die sich auf die unendlich vielen Werte des agM. bezieht, ist uns nicht erhalten. Dagegen steht der »nexus mutuus« auf einem Zettel, der oben S. 218 ff., art. [2.]—[6.], abgedruckt ist, und der schon durch das Wasserzeichen FHF 1810 auf eine wesentlich spätere Zeit verwiesen wird; er dürfte (vergl. die Bemerkungen oben S. 282) etwa 1825 geschrieben sein. Die hier in Betracht kommende Stelle ist namentlich die Formel
\begin{equation*}
\tag*{(*)}\frac1{(\mu)}=\frac1\mu+\frac{4ik}\lambda,
\end{equation*}
\noindent wo
\[
\mu=M(m,n),\quad\lambda=M(m,\sqrt{m^2-n^2})
\]
\noindent ist und \((\mu)\) einen der Werte bedeutet, die man für das agM. zwischen \(m\) und \(n\) erhält, wenn man »für ein \(n',n'',n'''\) etc. einen negativen Wert wählt«. Die Formel (*) setzt überdies voraus, daß wenn der reale Teil von \(\frac\mu\lambda\) wesentlich positiv ist, das gleiche auch für den realen Teil von \(\frac{(\mu)}\lambda\) gilt (vergl. oben S. 281). Der Weg, auf dem \textsc{Gauss} zu diesem Ergebnis vorgedrungen ist, kann, da Aufzeichnungen aus der Zeit der Entdeckung fehlen, nicht mit voller Sicherheit angegeben werden. Die Aufzeichnung von 1825 weist darauf hin, daß \textsc{Gauss} die Formel (*) wohl aus der Darstellung der \(m,n,\sqrt{m^2-n^2}\) mit Hilfe der Thetanullreihen \(p,q,r\) gefunden hat (siehe besonders den art. [6.] oben S. 222, 223, wo der allgemeinste Wert von \(M\) aufgestellt wird, für den \(\frac{q^2}{p^2}=\frac nm\) ungeändert bleibt). Die Wendungen »ex asse perspiciendi« und das sonst nur bei zahlentheoretischen Ergebnissen vorkommende »felicitas nobis est facta« (siehe die Nummern 30, 73, 114, 141 des \textit{Tagebuchs}) geben der Annahme Raum, daß \textsc{Gauss} auch sogleich die Beziehungen zur Theorie der binären quadratischen Formen durchschaut haben wird.
\par\indent Das \textls[120]{harmonische Mittel} zwischen zwei Zahlen \(m,n\) ist durch die Gleichung
\[
m_1=\frac{2mn}{m+n}
\]
\noindent erklärt. Der Algorithmus des \textls[120]{harmonisch-geometrischen Mittels} ist demnach:
\[
\begin{gathered}
\frac1{m_{\varkappa+1}}=\frac12\left(\frac1{m_\varkappa}+\frac1{n_\varkappa}\right),\quad\frac1{n_{\varkappa+1}}=\sqrt{\frac1{m_\varkappa}\cdot\frac1{n_\varkappa}}\\
(\varkappa=0,1,2,\ldots;\ m_0=m,\ n_0=n),
\end{gathered}
\]
\noindent so daß sich
\[
\lim_{\varkappa\to\infty}m_\varkappa=\lim_{\varkappa\to\infty}n_\varkappa=\frac{mn}{M(a,b)}
\]
\noindent ergibt, wodurch das harmonisch-geometrische Mittel auf das agM. zurückgeführt ist\footnote{Vergl. \textsc{Th. Lohnstein}, Zeitschrift für Math. und Physik \textbf{33} (1888), S. 316.}. In der Scheda Ac findet sich auf S. 20 die folgende Notiz:
\par\endgroup

\clearpage\SourceLeaf{580}\ArithmeticHead{551}{R}{1800 IUNIO 3—1800 IUN. 12.}
\indent Terminus constans expressionis
\[
\frac{A\,d\varphi}{\surd(f+2g\cos\varphi+h\cos\varphi^2)}
\]
\noindent est Medium Geometrico harmonicum inter
\[
\frac{A}{\mathord{\raisebox{-.25ex}{\scalebox{1.45}[2.0]{$\surd$}}}\dfrac{\surd((f+h)^2-4gg)+f-h}{2}}
\quad\text{et}\quad\frac{A}{{}^{4}\!\surd((f+h)^2-4gg)}.
\]
\medskip\begingroup\small
\noindent An einer andern, oben S. 14—16 abgedruckten Stelle der Scheda Ac, hat \textsc{Gauss} in Verbindung mit gewissen mittleren Werten der Zahlentheorie ein \textls[120]{arithmetisch-harmonisches Mittel} betrachtet (siehe besonders S. 16, art. [3.]).
\par\noindent\hfill\textsc{Klein.\quad Schlesinger.}
\par\endgroup
